\documentclass[11pt]{article}
\usepackage[margin=1in]{geometry}
\usepackage{enumitem}
% =============================================================================
% Preambles/header.tex — shared LaTeX/Beamer preamble for this project
%
% Use from a Beamer lecture by prepending:
%     \input{header}
% and compiling with TEXINPUTS=../Preambles:$TEXINPUTS (the /compile-latex
% skill does this automatically).
%
% PALETTE CONTRACT. Color names defined here MUST match the names in
% Quarto/theme-template.scss so Beamer and Quarto renderings use the same
% palette. The scripts/check-palette-sync.sh script enforces this.
%
% To customize: edit the HEX values below AND the matching $variable in
% Quarto/theme-template.scss, then run scripts/check-palette-sync.sh.
% =============================================================================

% -----------------------------------------------------------------------------
% Engine detection + encoding
%
% Our /compile-latex skill uses XeLaTeX, where inputenc is unnecessary and
% can produce encoding warnings. Only load inputenc under pdfTeX.
% -----------------------------------------------------------------------------
\usepackage{iftex}
\ifPDFTeX
  \usepackage[utf8]{inputenc}
\fi

\usepackage{amsmath, amssymb, amsthm}
\usepackage{booktabs}
\usepackage{graphicx}

% -----------------------------------------------------------------------------
% PALETTE — mirrors Quarto/theme-template.scss variable names.
% Load xcolor and define colors BEFORE anything that references them
% (notably hyperref, hypersetup, and the Beamer color assignments).
% -----------------------------------------------------------------------------
\usepackage{xcolor}

\definecolor{primary-blue}{HTML}{012169}      % headings, accents
\definecolor{primary-gold}{HTML}{B9975B}      % emphasis, borders
\definecolor{highlight-yellow}{HTML}{F2A900}  % markers, alerts
\definecolor{light-bg}{HTML}{E8EDF5}          % title / section backgrounds
\definecolor{jet}{HTML}{1A1A1A}               % body text

% Semantic colors (used in diagrams and annotations)
\definecolor{positive}{HTML}{15803D}          % good / observed / identified
\definecolor{negative}{HTML}{B91C1C}          % bad / problematic / bias
\definecolor{neutral}{HTML}{525252}           % reference / context

% Highlight accents (match .hi-* classes in the SCSS)
\definecolor{hi-slate}{HTML}{314F4F}
\definecolor{hi-green}{HTML}{2E8B57}
\definecolor{hi-red}{HTML}{C41E3A}

% Semantic aliases so skill templates and snippets can write
% \textcolor{accent}{...} without hard-coding "primary-blue".
\colorlet{accent}{primary-blue}
\colorlet{accent2}{primary-gold}

% -----------------------------------------------------------------------------
% Hyperref — loaded AFTER xcolor + palette so link colors resolve.
% -----------------------------------------------------------------------------
\usepackage{hyperref}
\hypersetup{colorlinks=true, linkcolor=primary-blue, urlcolor=primary-blue,
            citecolor=primary-blue, pdfborder={0 0 0}}

% -----------------------------------------------------------------------------
% Course code — set once per deck (after \input{header}, before \begin{document})
% via \coursecode{CS401}. Shown in the footer on every slide so a deck's course
% is identifiable without opening the title page. Empty by default.
% -----------------------------------------------------------------------------
\makeatletter
\newcommand{\coursecode}[1]{\gdef\@coursecodetext{#1}}
\gdef\@coursecodetext{}
\makeatother

% -----------------------------------------------------------------------------
% Beamer theme assignments (only apply under Beamer).
%
% We detect Beamer via \@ifclassloaded{beamer}, which is the documented,
% non-internal way. This must be wrapped in \makeatletter/\makeatother
% because \@ifclassloaded contains an @-catcode character.
% -----------------------------------------------------------------------------
\makeatletter
\@ifclassloaded{beamer}{%
  \usefonttheme{professionalfonts}
  \setbeamercolor{structure}{fg=primary-blue}
  \setbeamercolor{frametitle}{fg=primary-blue}
  \setbeamercolor{title}{fg=primary-blue}
  \setbeamercolor{subtitle}{fg=primary-gold}
  \setbeamercolor{author}{fg=primary-blue}
  \setbeamercolor{institute}{fg=neutral}
  \setbeamercolor{date}{fg=primary-gold}
  \setbeamercolor{itemize item}{fg=primary-blue}
  \setbeamercolor{itemize subitem}{fg=primary-gold}
  \setbeamercolor{alerted text}{fg=primary-gold}
  \setbeamercolor{block title}{fg=white, bg=primary-blue}
  \setbeamercolor{block body}{bg=light-bg}
  % Semantic block variants: block=definition/concept (blue), exampleblock=
  % worked example (green/positive), alertblock=key takeaway or caution (gold).
  % Keep to <=2 colored boxes per slide across ALL three variants (INV-7).
  \setbeamercolor{block title example}{fg=white, bg=positive}
  \setbeamercolor{block body example}{bg=light-bg}
  \setbeamercolor{block title alerted}{fg=white, bg=primary-gold}
  \setbeamercolor{block body alerted}{bg=light-bg}

  % Minimal footer: course code (left) + page X / N (right), no navigation clutter
  \setbeamertemplate{navigation symbols}{}
  \setbeamertemplate{footline}{%
    \color{neutral}\footnotesize\hspace{0.7em}\@coursecodetext\hfill
    \insertframenumber/\inserttotalframenumber\hspace{0.7em}\vspace{0.3em}}
}{}
\makeatother

% -----------------------------------------------------------------------------
% TikZ setup — libraries the snippet gallery uses
% -----------------------------------------------------------------------------
\usepackage{tikz}
\usetikzlibrary{arrows.meta, positioning, calc, decorations.pathreplacing,
                fit, shapes.geometric, backgrounds}

% Shared TikZ styles — snippets and hand-written diagrams can pick these up
% via `\begin{tikzpicture}[every node/.style={...}]` or by naming specific
% styles (e.g., `\node[dag-node] (X) at (0,0) {X};`).
\tikzset{
  % Node styles for causal DAGs and similar
  dag-node/.style = {
    draw=primary-blue, fill=light-bg, rounded corners=2pt,
    minimum width=1.2cm, minimum height=0.9cm, align=center, font=\small
  },
  decision-node/.style = {
    draw=primary-gold, fill=white, diamond, aspect=1.8,
    minimum width=1.8cm, minimum height=1.2cm, align=center, font=\small,
    inner sep=1pt
  },
  % Edge styles — observed vs counterfactual
  observed-edge/.style = {-{Latex[length=2.5mm]}, thick, draw=jet},
  counterfactual-edge/.style = {-{Latex[length=2.5mm]}, thick, draw=neutral, dashed},
  confound-edge/.style = {-{Latex[length=2.5mm]}, thick, draw=negative, dashed},
  % Data-point styles for plots
  observed-dot/.style = {fill=primary-blue, draw=primary-blue, circle, inner sep=1.2pt},
  counterfactual-dot/.style = {fill=white, draw=neutral, circle, inner sep=1.2pt}
}

% -----------------------------------------------------------------------------
% Convenience macros
% -----------------------------------------------------------------------------
\newcommand{\muted}[1]{\textcolor{neutral}{#1}}
\newcommand{\key}[1]{\textcolor{primary-gold}{\textbf{#1}}}
\newcommand{\good}[1]{\textcolor{positive}{#1}}
\newcommand{\bad}[1]{\textcolor{negative}{#1}}

% Standout frame macro: full-bleed dark background for section transitions.
% Beamer-only (uses \begin{frame}/\setbeamercolor/\usebeamerfont) -- do not
% call from an article-class file (e.g. Notes/<CODE>/*-notes.tex). \muted,
% \key, \good, \bad, and \coursecode above are plain xcolor/gdef and are
% safe to use from either class; verified when Notes/article-class support
% was added.
% Expand: \transitionslide{Your transition title}
\newcommand{\transitionslide}[1]{%
  \begingroup
  \setbeamercolor{background canvas}{bg=primary-blue}
  \begin{frame}[plain]
    \centering
    \vfill
    {\usebeamerfont{title}\color{white}\Large #1\par}
    \vfill
  \end{frame}
  \endgroup
}

% Section divider frame (Beamer-only), an alternative to \transitionslide with a
% darker two-line style: near-black (jet) background, a small white label line
% above a larger gold title, and a thin gold rule along the bottom edge.
% Expand: \sectiondivider{Section 2}{Pointers}
\newcommand{\sectiondivider}[2]{%
  \begingroup
  \setbeamercolor{background canvas}{bg=jet}
  \begin{frame}[plain]
    \vspace*{3.0cm}
    {\color{white}\ttfamily\large #1\par}
    \vspace{0.5cm}
    {\color{primary-gold}\LARGE #2\par}
    \begin{tikzpicture}[remember picture, overlay]
      \draw[primary-gold, line width=2.5pt]
        ($(current page.south west)+(0,0.1)$) -- ($(current page.south east)+(0,0.1)$);
    \end{tikzpicture}
  \end{frame}
  \endgroup
}

% -----------------------------------------------------------------------------
% Channel watermark — "Concepts That Click" (YouTube).
%
% Article-class files (CompetitiveExam Q/A sets, Notes, Assignments) get a
% light diagonal draft watermark on every page plus a centered footer naming
% the channel. Beamer decks get a subtle TikZ background watermark instead
% (the footline already carries the course code). Centralized here so every
% MCQ PDF is branded without per-file edits.
% -----------------------------------------------------------------------------
\makeatletter
\@ifclassloaded{article}{%
  \usepackage{draftwatermark}
  \SetWatermarkText{Concepts That Click}
  \SetWatermarkScale{1}
  \SetWatermarkFontSize{2cm}
  \SetWatermarkLightness{0.86}
  \SetWatermarkAngle{45}
  \usepackage{fancyhdr}
  \pagestyle{fancy}
  \fancyhf{}
  \fancyfoot[C]{\footnotesize\textcolor{neutral}{Concepts That Click~|~YouTube: \href{https://www.youtube.com/@concepts.thatclick}{@concepts.thatclick}}}
  \renewcommand{\headrulewidth}{0pt}
  \renewcommand{\footrulewidth}{0pt}
  \fancypagestyle{plain}{%
    \fancyhf{}%
    \fancyfoot[C]{\footnotesize\textcolor{neutral}{Concepts That Click~|~YouTube: \href{https://www.youtube.com/@concepts.thatclick}{@concepts.thatclick}}}%
    \renewcommand{\headrulewidth}{0pt}%
    \renewcommand{\footrulewidth}{0pt}%
  }
}{}
\@ifclassloaded{beamer}{%
  \setbeamertemplate{background}{%
    \begin{tikzpicture}[remember picture, overlay]
      \node[rotate=45, opacity=0.055, align=center]
        at (current page.center)
        {\Huge\textcolor{neutral}{Concepts That Click}};
    \end{tikzpicture}%
  }
}{}
\makeatother 
\coursecode{JKSSB}
\title{Lesson 56 Practice: Cybersecurity, IT Act and Security Controls}
\author{JKSSB Computer Awareness}
\date{\today}
\begin{document}
\maketitle
\noindent\textbf{Provenance:} All 20 questions are original JKSSB-pattern
questions (\textit{[Original]}); none reproduces an official paper item.
Brand names appear only as distractors, never as examined facts.

\begin{enumerate}[label=\textbf{Q\arabic*.}, leftmargin=*]
\item \textit{[Original]} Expand MFA.
  \begin{enumerate}[label=\Alph*.]
    \item Multi-Factor Authentication
    \item Multiple-File Access
    \item Managed Firewall Authority
    \item Master File Archive
  \end{enumerate}
\item \textit{[Original]} SSL stands for:
  \begin{enumerate}[label=\Alph*.]
    \item Secure Sockets Layer
    \item System Security Lock
    \item Shared Signature Log
    \item Synchronised Server Link
  \end{enumerate}
\item \textit{[Original]} The Information Technology Act of India was
  enacted in the year:
  \begin{enumerate}[label=\Alph*.]
    \item 1998
    \item 2000
    \item 2008
    \item 2013
  \end{enumerate}
\item \textit{[Original]} Which of the following is \textbf{NOT} within the
  scope of the Information Technology Act, 2000 as covered in this lesson?
  \begin{enumerate}[label=\Alph*.]
    \item Legal validity of electronic records and electronic signatures
    \item Electronic governance
    \item Cyber contraventions and offences
    \item Regulation of printer hardware prices and typing-test fees
  \end{enumerate}
\item \textit{[Original]} Which term is the \textbf{umbrella} term for
  malicious software, under which virus, worm, Trojan, and ransomware fall?
  \begin{enumerate}[label=\Alph*.]
    \item Malware
    \item Firewall
    \item Patch
    \item Backup
  \end{enumerate}
\item \textit{[Original]} Which malware attaches itself to a host file and
  needs human action, such as opening an infected file, to spread?
  \begin{enumerate}[label=\Alph*.]
    \item Worm
    \item Trojan horse
    \item Virus
    \item Ransomware
  \end{enumerate}
\item \textit{[Original]} Which malware copies itself and spreads across a
  network without a host file and without waiting for the user to act?
  \begin{enumerate}[label=\Alph*.]
    \item Virus
    \item Worm
    \item Trojan horse
    \item Spyware-free utility
  \end{enumerate}
\item \textit{[Original]} Which statement is \textbf{NOT} true of a Trojan
  horse?
  \begin{enumerate}[label=\Alph*.]
    \item It pretends to be useful software while hiding a harmful purpose.
    \item It opens a backdoor for stealing passwords or controlling the machine.
    \item It replicates itself across networks without any user action.
    \item It relies on the user installing it under a false disguise.
  \end{enumerate}
\item \textit{[Original]} Ransomware is best described as malware that:
  \begin{enumerate}[label=\Alph*.]
    \item filters network traffic against a set of rules.
    \item encrypts or locks the victim's files and demands payment for release.
    \item isolates suspicious files for later inspection.
    \item fixes software flaws by installing vendor updates.
  \end{enumerate}
\item \textit{[Original]} Evaluate the statements:
  \begin{enumerate}[label=\Roman*.]
    \item A virus needs a host file and human action to spread.
    \item A worm is standalone and spreads across a network unaided.
    \item A Trojan horse replicates itself automatically like a worm.
  \end{enumerate}
  \begin{enumerate}[label=\Alph*.]
    \item I only
    \item I and II only
    \item II and III only
    \item I, II, and III
  \end{enumerate}
\item \textit{[Original]} Placing a suspicious file in quarantine means:
  \begin{enumerate}[label=\Alph*.]
    \item deleting it permanently beyond recovery.
    \item isolating it so it cannot run or spread, pending a decision.
    \item backing it up to a separate disk for safekeeping.
    \item encrypting it for secure transmission to another user.
  \end{enumerate}
\item \textit{[Original]} Phishing is:
  \begin{enumerate}[label=\Alph*.]
    \item a fake message impersonating a trusted sender to steal passwords or money.
    \item a barrier that filters incoming and outgoing network traffic.
    \item a vendor update that fixes known security flaws.
    \item a protocol that encrypts data during transmission.
  \end{enumerate}
\item \textit{[Original]} A customer receives a mail reading ``your account
  is blocked; click here and enter your OTP to restore it.'' The correct
  anti-phishing response is to:
  \begin{enumerate}[label=\Alph*.]
    \item click at once and enter the OTP before the deadline.
    \item forward the OTP to the sender for verification.
    \item ignore the link, share no OTP, and verify through the known official channel.
    \item download the attached ``security tool'' and run it.
  \end{enumerate}
\item \textit{[Original]} Encryption primarily provides:
  \begin{enumerate}[label=\Alph*.]
    \item confidentiality, by converting readable data so only authorised persons can read it.
    \item proof of the sender's identity without hiding the message.
    \item filtering of network packets at the boundary.
    \item isolation of suspicious files from the system.
  \end{enumerate}
\item \textit{[Original]} A digital signature primarily provides:
  \begin{enumerate}[label=\Alph*.]
    \item secrecy of the message content from eavesdroppers.
    \item removal of malware already present on the machine.
    \item proof of who sent the message and that it was not altered.
    \item an encrypted channel for data in transit.
  \end{enumerate}
\item \textit{[Original]} SSL/TLS protect:
  \begin{enumerate}[label=\Alph*.]
    \item files already infected on the disk, by cleaning the malware.
    \item data in transit, by providing a secure encrypted channel.
    \item login accounts, by demanding an OTP at every step.
    \item software flaws, by installing vendor-supplied updates.
  \end{enumerate}
\item \textit{[Original]} Which control filters incoming and outgoing
  network traffic against a set of rules and blocks unauthorised access?
  \begin{enumerate}[label=\Alph*.]
    \item Firewall
    \item Norton
    \item McAfee
    \item Quick Heal
  \end{enumerate}
\item \textit{[Original]} Patching means:
  \begin{enumerate}[label=\Alph*.]
    \item demanding payment to release locked files.
    \item installing vendor-supplied updates that fix known security flaws.
    \item adding a second login check such as an OTP.
    \item copying data to a separate store for recovery after loss.
  \end{enumerate}
\item \textit{[Original]} An attacker obtains a user's password, but the
  account also requires an OTP sent to the user's phone, which the attacker
  does not have. The login fails. This scenario shows the value of:
  \begin{enumerate}[label=\Alph*.]
    \item quarantine
    \item ransomware
    \item Multi-Factor Authentication (MFA)
    \item phishing
  \end{enumerate}
\item \textit{[Original]} Which pairing is correctly matched?
  \begin{enumerate}[label=\Alph*.]
    \item Quarantine --- deletes the file permanently beyond recovery
    \item Kaspersky --- a network barrier that filters traffic by rules
    \item Encryption --- proves the sender without hiding the message
    \item Patching --- fixes security flaws by installing updates
  \end{enumerate}
\end{enumerate}
\end{document}
